\documentclass{amsart}
\usepackage{amsmath,amssymb,amsthm,mathtools,enumitem}
\title{Homework 5}
\author{Antony Perez}
\begin{document}
\maketitle
\subsection*{4.1}
Find the inverse of the matrices
\[
\begin{pmatrix}
1 & 2 & 1\\
3 & 7 & 3\\
2 & 3 & 4
\end{pmatrix},
\qquad
\begin{pmatrix}
1 & -1 & 2\\
1 & 1 & -2\\
1 & 1 & 4
\end{pmatrix}.
\]

Show all steps.
\begin{proof}
  $R_2 \to R_2 - 3R_1;
R_3 \to R_3 - 2R_1;
R_3 \to R_3 + R_2;
R_3 \to \tfrac{1}{2}R_3;
R_1 \to R_1 - R_3;
R_1 \to R_1 - 2R_2$;
\[
A^{-1} = \begin{pmatrix} \tfrac{19}{2} & -\tfrac{5}{2} & -\tfrac{1}{2} \\ -3 & 1 & 0 \\ -\tfrac{5}{2} & \tfrac{1}{2} & \tfrac{1}{2} \end{pmatrix}
\]

$R_2 \to R_2 - R_1;
R_3 \to R_3 - R_1;
R_3 \to R_3 - R_2;
R_2 \to \tfrac{1}{2}R_2;
R_3 \to \tfrac{1}{6}R_3;
R_1 \to R_1 + R_2;
R_2 \to R_2 + 2R_3;$

\[A^{-1} = \begin{pmatrix} \tfrac{1}{2} & \tfrac{1}{2} & 0 \\ -\tfrac{1}{2} & \tfrac{1}{6} & \tfrac{1}{3} \\ 0 & -\tfrac{1}{6} & \tfrac{1}{6} \end{pmatrix}\]
\end{proof}
\subsection*{5.1}
True or false:
\begin{enumerate}[label=\alph*)]
  \item Every vector space that is generated by a finite set has a basis;
  \item Every vector space has a (finite) basis;
  \item A vector space cannot have more than one basis;
  \item If a vector space has a finite basis, then the number of vectors in every basis is the same.
  \item The dimension of $\mathbb{P}_n$ is $n$;
  \item The dimension on $M_{m\times n}$ is $m+n$;
  \item If vectors $\mathbf{v}_1, \mathbf{v}_2, \ldots, \mathbf{v}_n$ generate (span) the vector space $V$, then every vector in $V$ can be written as a linear combination of vector $\mathbf{v}_1, \mathbf{v}_2, \ldots, \mathbf{v}_n$ in only one way;
  \item Every subspace of a finite-dimensional space is finite-dimensional;
  \item If $V$ is a vector space having dimension $n$, then $V$ has exactly one subspace of dimension $0$ and exactly one subspace of dimension $n$.
\end{enumerate}
\begin{proof}
    TFFTFFFTT
\end{proof}
\subsection*{5.2}
Prove that if $V$ is a vector space having dimension $n$, then a system of vectors $\mathbf{v}_1, \mathbf{v}_2, \ldots, \mathbf{v}_n$ in $V$ is linearly independent if and only if it spans $V$.
\begin{proof}
  Suppose $\mathbf{v}_1,\ldots,\mathbf{v}_n$ are linearly independent, but do not span $V$. Thus, there exists $\mathbf{w} \notin \text{span} \{\mathbf{v}_1,\ldots,\mathbf{v}_n\}$. Thus, $\mathbf{w}$ is linearly independent. Now, we may write $a_1\mathbf{v}_1+\ldots+a_2\mathbf{v}_n+b\mathbf{w}$. Thus, by definition, $b=0$. If it werent, we'd get $\mathbf{w} = -\frac{1}{b}(a_1\mathbf{v}_1 + \ldots + a_n \mathbf{v}_n)$, which contradicts the fact that $\mathbf{w}$ is independent. Thus, we have $n+1$ vectors now. But there can be at most $n$ independent vectors in $V$. Thus, $v_1,\ldots, v_n$ span $V$.

  Suppose $\mathbf{v}_1,\ldots,\mathbf{v}_n$ span $V$, but are not linearly independent. Thus, there is some vector $v_i$ that we may remove from the span, giving us $n-1$ vectors. But in order to span $V$, the spanning set must have at least $n$ vectors. Thus, $\mathbf{v}_1, \ldots, \mathbf{v}_n$ are linearly independent. 
\end{proof}
\subsection*{5.3}
Prove that a linearly independent system of vectors $\mathbf{v}_1, \mathbf{v}_2, \ldots, \mathbf{v}_n$ in a vector space $V$ is a basis if and only if $n = \dim V$.
\begin{proof}
    We are given $\mathbf{v}_1, \mathbf{v}_2, \ldots, \mathbf{v}_n$ linearly independent vectors in $V$. Now say that this system spans $V$. Thus, this is a basis for $V$. By definition, $\dim V = n$, because there are exactly $n$ vectors in our basis that generate $V$. 
    \noindent
    Conversely, assume $\dim V = n$. We must show that our linearly independent vectors span $V$. Hence, suppose there exists $\mathbf{w}\in V$ such that $\mathbf{w} \not \in \text{span}\{\mathbf{v}_1, \mathbf{v}_2, \ldots, \mathbf{v}_n\}$. Suppose $a_1\mathbf{v}_1+ \mathbf{v}_2+ \ldots+ a_n\mathbf{v}_n + b \mathbf{w}=0$. Assume $b \neq 0$. We get $\mathbf{w} = -\frac{1}{b}\sum_{i=1}^{n}a_i\mathbf{v}_i$. But we said $\mathbf{w}$ isn't in the span. Thus, $b=0$. This shows $\mathbf{w}$ is linearly independent. Thus, we get $a_1\mathbf{v}_1+ a_2\mathbf{v}_2+ \ldots+ a_n\mathbf{v}_n + b\mathbf{w}=0.$ Now, we have $n+1$ linearly independent vectors. But $\dim V =n$, and there can be at most $n$ linearly independent vectors in $V$. This is a contradiction. Therefore, our original system $\{\mathbf{v}_1, \mathbf{v}_2, \ldots, \mathbf{v}_n\}$ spans $V$, making it a basis for $V$. 
\end{proof}

\subsection*{5.4}
(An old problem revisited: now this problem should be easy) Is it possible that vectors $\mathbf{v}_1, \mathbf{v}_2, \mathbf{v}_3$ are linearly dependent, but the vectors $\mathbf{w}_1 = \mathbf{v}_1 + \mathbf{v}_2$, $\mathbf{w}_2 = \mathbf{v}_2 + \mathbf{v}_3$ and $\mathbf{w}_3 = \mathbf{v}_3 + \mathbf{v}_1$ are linearly \emph{independent}? Hint: What dimension can the subspace $\operatorname{span}(\mathbf{v}_1, \mathbf{v}_2, \mathbf{v}_3)$ have?
\begin{proof}
    No. Since $\mathbf{v}_1, \mathbf{v}_2, \mathbf{v}_3$ are linearly dependent, they are linear combinations of the others, so $\dim (\mathbf{v}_1, \mathbf{v}_2, \mathbf{v}_3) \le 2$. Since $\mathbf{w}_1, \mathbf{w}_2, \mathbf{w}_3$ are linear combinations of $\mathbf{v}_1, \mathbf{v}_2, \mathbf{v}_3$, they lie in $\operatorname{span}(\mathbf{v}_1, \mathbf{v}_2, \mathbf{v}_3)$. Since there are three vectors in an, at most, 2 dimensional space, one of them must be linearly dependent. Therefore, $\mathbf{w}_1, \mathbf{w}_2, \mathbf{w}_3$ cannot be linearly independent.
\end{proof}

\subsection*{5.5}
Let vectors $\mathbf{u}, \mathbf{v}, \mathbf{w}$ be a basis in $V$. Show that $\mathbf{u} + \mathbf{v} + \mathbf{w}$, $\mathbf{v} + \mathbf{w}$, $\mathbf{w}$ is also a basis in $V$.

\begin{proof}
Suppose $a(u+v+w) + b(v+w) + cw = 0$. Then $au + (a+b)v + (a+b+c)w = 0$, so by linear independence of $u,v,w$, $a = a+b = a+b+c = 0$. Hence $u+v+w,\, v+w,\, w$ are linearly independent. They also span $V$, since $u = (u+v+w)-(v+w)$, $v = (v+w)-w$, $w = w$. Thus they form a basis.
\end{proof}

\subsection*{5.6}
Consider in the space $\mathbb{R}^5$ vectors $\mathbf{v}_1 = (2, -1, 1, 5, -3)^T$, $\mathbf{v}_2 = (3, -2, 0, 0, 0)^T$, $\mathbf{v}_3 = (1, 1, 50, -921, 0)^T$.
\begin{enumerate}[label=\alph*)]
  \item Prove that these vectors are linearly independent.
  \item Complete this system of vectors to a basis.
\end{enumerate}
If you do part b) first you can do everything without any computations.

\begin{proof}
  By the Rank Theorem, we can construct the matrix and row reduce:\\ $\begin{pmatrix}
    1 & 1 & 50 & -921 & 0 \\ 
    3 & -2 &0&0 & 0 \\ 
    2 & -1 & 1 & 5 & 3
  \end{pmatrix} \sim \begin{pmatrix}
    1 & 1 & 50 & -921 & 0 \\ 
    0 & -5 & -150 &  2763 & 0 \\ 
    0 & 0 & -45 & 946 & 15
  \end{pmatrix}$.
  Thus, by rank theorem, we see that the three vectors are linearly independent. Completing the basis, we get $\begin{pmatrix}
    1 & 1 & 50 & -921 & 0 \\ 
    0 & -5 & -150 &  2763 & 0 \\ 
    0 & 0 & -45 & 946 & 15 \\ 
    0 & 0 & 0 & 1 & 0 \\ 
   0 & 0 & 0 & 0 & 1 
  \end{pmatrix}$. Thus, the basis is \\$\{(2, -1, 1, 5, -3)^T,(3, -2, 0, 0, 0)^T,(1, 1, 50, -921, 0)^T, (0,0,0,1,0)^T, (0,0,0,0,1)^T\}$
\end{proof}
\subsection*{6.1}
\begin{enumerate}[label=\alph*)]
  \item Any system of linear equations has at least one solution;
  \item Any system of linear equations has at most one solution;
  \item Any homogeneous system of linear equations has at least one solution;
  \item Any system of $n$ linear equations in $n$ unknowns has at least one solution;
  \item Any system of $n$ linear equations in $n$ unknowns has at most one solution;
  \item If the homogeneous system corresponding to a given system of a linear equations has a solution, then the given system has a solution;
  \item If the coefficient matrix of a homogeneous system of $n$ linear equations in $n$ unknowns is invertible, then the system has no non-zero solution;
  \item The solution set of any system of $m$ equations in $n$ unknowns is a subspace in $\mathbb{R}^n$;
  \item The solution set of any homogeneous system of $m$ equations in $n$ unknowns is a subspace in $\mathbb{R}^n$.
\end{enumerate}
\begin{proof}
  FFTFFFTFT
\end{proof}

\subsection*{6.2}
Find a $2 \times 3$ system (2 equations with 3 unknowns) such that its general solution has a form
\[
\begin{pmatrix} 1 \\ 1 \\ 0 \end{pmatrix}
+ s \begin{pmatrix} 1 \\ 2 \\ 1 \end{pmatrix}, \qquad s \in \mathbb{R}.
\]
\begin{proof}
    $\begin{bmatrix}
        1 & 0 & -1 \\ 
        0 & 1 & -2 \\ 
    \end{bmatrix}\textbf{x} =\begin{bmatrix}
        1 \\ 1
    \end{bmatrix}$. 
\end{proof}
\subsection*{7.1}
\begin{enumerate}[label=\alph*)]
  \item The rank of a matrix is equal to the number of its non-zero columns;
  \item The $m \times n$ zero matrix is the only $m \times n$ matrix having rank $0$;
  \item Elementary row operations preserve rank;
  \item Elementary column operations do not necessarily preserve rank;
  \item The rank of a matrix is equal to the maximum number of linearly independent columns in the matrix;
  \item The rank of a matrix is equal to the maximum number of linearly independent rows in the matrix;
  \item The rank of an $n \times n$ matrix is at most $n$;
  \item An $n \times n$ matrix having rank $n$ is invertible.
\end{enumerate}
\begin{proof}
  FTTFTTTT
\end{proof}
\subsection*{7.2}
A $54 \times 37$ matrix has rank $31$. What are dimensions of all 4 fundamental subspaces?
\begin{proof}
    Column Space: 31 \\
    Null Space: 6 \\
    Row Space: 31 \\ 
    Left Null Space: 23
\end{proof}
\subsection*{7.3}
Compute rank and find bases of all four fundamental subspaces for the matrices
\[
\begin{pmatrix}
1 & 1 & 0\\
0 & 1 & 1\\
1 & 1 & 0
\end{pmatrix},
\qquad
\begin{pmatrix}
1 & 2 & 3 & 1 & 1\\
1 & 4 & 0 & 1 & 2\\
0 & 2 & -3 & 0 & 1\\
1 & 0 & 0 & 0 & 0
\end{pmatrix}.
\]

\begin{proof}
    For first matrix: \\ 
    Column Space: 2 \\ 
    Null space: 1 \\ 
    Row Space: 2 \\ 
    Left null space: 1 \\ 
   $\operatorname{Ran}A$: $(1,0,1)^T,\ (1,1,1)^T$ (pivot columns of $A$).\\
   $\operatorname{Ran}A^T$: $(1,0,-1),\ (0,1,1)$.\\
   $\operatorname{Ker}A$: $(1,-1,1)^T$.\\
   $\operatorname{Ker}A^T$: $(1,0,-1)^T$. \\
Second matrix: 

   $\operatorname{Ran}B$: $(1,1,0,1)^T,\ (2,4,2,0)^T,\ (3,0,-3,0)^T$ (pivot columns of $B$).\\
   $\operatorname{Ran}B^T$: $(1,0,0,0,0),\ (0,2,3,1,1),\ (0,0,-6,-1,0)$.\\
   $\operatorname{Ker}B$: $(0,-3,-2,12,0)^T,\ (0,-1,0,0,2)^T$.\\
   $\operatorname{Ker}B^T$: $(1,-1,1,0)^T$ (since row$_3$ = row$_2$ $-$ row$_1$).\\

    Col Space: 3 \\ 
    Null space: 2 \\ 
    Row Space: 3 \\ 
    Left Null Space : 1
\end{proof}
\subsection*{7.4}
Prove that if $A : X \to Y$ and $V$ is a subspace of $X$ then $\dim AV \le \operatorname{rank} A$. ($AV$ here means the subspace $V$ transformed by the transformation $A$, i.e.\ any vector in $AV$ can be represented as $A\mathbf{v}$, $\mathbf{v} \in V$). Deduce from here that $\operatorname{rank}(AB) \le \operatorname{rank} A$.

Remark: Here one can use the fact that if $V \subset W$ then $\dim V \le \dim W$. Do you understand why is it true?

\begin{proof}
$AV$ is a subspace: $\mathbf{0} = A\mathbf{0} \in AV$, and $\alpha A\mathbf{v}_1 + \beta A\mathbf{v}_2 = A(\alpha\mathbf{v}_1 + \beta\mathbf{v}_2) \in AV$ since $\alpha\mathbf{v}_1+\beta\mathbf{v}_2 \in V$. If $\mathbf{w} \in AV$, then $\mathbf{w} = A\mathbf{v}$ with $\mathbf{v}\in V\subset X$, so $\mathbf{w}\in\operatorname{Ran}A$. Thus $AV \subset \operatorname{Ran}A$, and $\dim AV \le \dim\operatorname{Ran}A = \operatorname{rank}A$.

For the second claim, $\operatorname{Ran}(AB) = A(\operatorname{Ran}B)$, so taking $V = \operatorname{Ran}B$ gives $\operatorname{rank}(AB) = \dim A(\operatorname{Ran}B) \le \operatorname{rank}A$.
\end{proof}
\subsection*{7.5}
Prove that if $A : X \to Y$ and $V$ is a subspace of $X$ then $\dim AV \le \dim V$. Deduce from here that $\operatorname{rank}(AB) \le \operatorname{rank} B$.
\begin{proof}
  Since $V$ is a v.s., dim$(V)=n \implies n$ vectors needed in basis. Let $v_1,\ldots,v_n$ be a basis for $V$. Thus, any $v \in V$ may be represented as $v = \sum_{i=1}^{n}\alpha_i v_i$. For any $w \in AV$, there exists some $v$ such that $w = Av$. Thus, $w = \sum_{i=1}^{n}\alpha_i A(v_i)$. Thus, it takes $n$ vectors to span $AV$, which shows dim$(AV) \le \text{dim}(V)$. Doing the same process, we arrive at the same result for $\text{Rank}(AB) \le \text{Rank}(B)$ when we let $V = \text{Ran}(B)$.  
\end{proof}

\subsection*{7.6}
Prove that if the product $AB$ of two $n \times n$ matrices is invertible, then both $A$ and $B$ are invertible. Even if you know about determinants, do not use them, we did not cover them yet. Hint: use previous 2 problems.
\begin{proof}
Since $AB$ is an invertible $n\times n$ matrix, $\operatorname{rank}(AB) = n$. By 7.4 and 7.5,
\[
n = \operatorname{rank}(AB) \le \operatorname{rank}A \le n, \qquad n = \operatorname{rank}(AB) \le \operatorname{rank}B \le n,
\]
so $\operatorname{rank}A = \operatorname{rank}B = n$. An $n\times n$ matrix of rank $n$ is invertible, so $A$ and $B$ are invertible.
\end{proof}

\subsection*{7.7}
Prove that if $A\mathbf{x} = \mathbf{0}$ has unique solution, then the equation $A^T\mathbf{x} = \mathbf{b}$ has a solution for every right side $\mathbf{b}$.

Hint: count pivots

\begin{proof}
Let $A$ be $m\times n$. Since $A\mathbf{x}=\mathbf{0}$ has a unique solution, there are no free variables, so the echelon form of $A$ has a pivot in every column. Hence $\operatorname{rank}A = n$, and so $\operatorname{rank}A^T = \operatorname{rank}A = n$. But $A^T$ is $n\times m$, so its $n$ pivots lie in $n$ different rows. Therefore $A^T\mathbf{x}=\mathbf{b}$ is consistent for every $\mathbf{b}$.
\end{proof}
\end{document}